% Formatted using the Springer LNCS samplepaper template style.
\documentclass[runningheads]{llncs}

\usepackage[T1]{fontenc}
\usepackage[utf8]{inputenc}
\usepackage{graphicx}
\usepackage{amsmath,amssymb}
\usepackage{url}
\usepackage{tabularx}
\usepackage{booktabs}
\usepackage{float}
\usepackage{array}

\begin{document}

\title{Belief Networks for Fault Diagnosis under Incomplete Information: Applicability and Limitations}
\titlerunning{Belief Networks for Fault Diagnosis}

\author{Pablo Garcia-Legaz Levi\inst{1}}
\authorrunning{P. Garcia-Legaz Levi}

\institute{Vrije Universiteit Amsterdam, Amsterdam, The Netherlands\\
Student number: 2727290}

\maketitle

\begin{abstract}
Fault diagnosis often requires reasoning from observable symptoms to hidden component failures under incomplete information. Bayesian networks are useful for structured diagnostic reasoning, but they require a complete specification of priors and conditional probability tables. This thesis investigates belief networks, represented using Dempster--Shafer belief functions, as a practical alternative when the diagnostic structure is known but some probabilistic information is missing or weakly justified. A modular electrical system is used as a case study. A complete Bayesian network provides a reference structure and validation point, while the Belief Function Machine in MATLAB is used to implement belief-network variants with progressively incomplete information. The results show that when evidence is decisive, the belief network preserves the correct diagnosis, while under missing priors or measurements it exposes unresolved alternatives through belief and plausibility rather than forcing precise posterior rankings. The thesis concludes with practical guidance on when belief networks are recommended for industrial fault diagnosis and where scalability, tooling, and interpretation remain limiting factors.

\keywords{Belief networks \and Dempster--Shafer theory \and Fault diagnosis \and Incomplete information \and Bayesian networks}
\end{abstract}

\section{Introduction}
\label{sec:introduction}

Industrial fault diagnosis often requires decisions before complete quantitative information is available. Engineers may understand the structure of a system, know which components can fail, and observe symptoms such as alarms, indicator lights, or sensor readings. However, they may not know all the numerical probabilities needed to quantify every possible fault-symptom relation. This thesis studies whether belief networks can support such diagnostic situations, and under what practical conditions they are useful for industrial users.

The diagnostic setting is introduced through a modular electrical system composed of one power supply module, eight control modules, and one load module. If the final lamp is off, several explanations may be possible: the battery may be exhausted, the power supply may be shorted, a switch may be detached, a cable may be disconnected, or the lamp itself may be faulty. Indicator LEDs and measurements can help localise the problem, but they are indirect evidence. A downstream indicator being off may indicate a local fault, but it may also be the consequence of an upstream interruption. The system is described in detail in Section~\ref{sec:case_study_diagnostic_model}; here it is used only to motivate the type of ambiguity that appears in fault diagnosis.

Bayesian networks are a widely used framework for reasoning under uncertainty in structured domains \cite{Pearl1988,Heckerman1995}. In fault diagnosis, a Bayesian network can represent hidden faults, internal system states, and observable symptoms as nodes in a directed acyclic graph. Evidence can then be propagated through the graph to update the posterior probabilities of possible faults. This is not in conflict with the motivation of the thesis: Bayesian networks do reason under uncertainty. The limitation considered here is more specific. A Bayesian network represents uncertainty through precise probability distributions, and therefore requires a complete probabilistic specification before standard inference can be performed.

This requirement becomes problematic when a diagnostic structure is known but some priors or conditional probability values are unavailable, weakly justified, or impossible to estimate reliably. In practice, a modeller may know that a detached switch interrupts voltage propagation, but not know the exact probability of each downstream symptom under every parent-state configuration. Completing the model with arbitrary assumptions, such as uniform probabilities or maximum-entropy estimates, can make inference possible, but may also produce precise posterior probabilities that hide the incompleteness of the underlying knowledge \cite{Shenoy2023IncompleteBN}.

This thesis investigates belief networks as a practical alternative for such cases. The term \emph{belief network} is used here to refer to a graphical diagnostic model represented using Dempster--Shafer belief functions. Unlike classical probability theory, belief-function theory allows support to be assigned not only to individual hypotheses, but also to sets of hypotheses \cite{Dempster1967,Shafer1976}. For example, if the evidence supports an upstream electrical interruption but does not justify distinguishing between a battery fault and an early cable fault, a belief function can assign support to the set containing both possibilities rather than forcing an artificial split between them.

The main outputs of belief-function inference are belief and plausibility. Belief represents support that is definitely committed to a proposition, while plausibility represents support that remains compatible with it \cite{Shafer1976,SmetsKennes1994}. In this thesis, the difference between plausibility and belief is not treated as an ordinary probability interval. Rather, it is interpreted as unresolved evidential support: the model may not have enough information to commit support to a specific fault, even though that fault remains compatible with the available evidence. This distinction is useful in diagnosis because it helps separate lack of support from evidence against a hypothesis.

The approach is motivated especially by Shenoy's work on inference in incomplete Bayesian networks using Dempster--Shafer belief functions \cite{Shenoy2023IncompleteBN}. Shenoy provides the conceptual bridge: when a Bayesian-style network structure is useful but some probabilities are missing, the known parts can be embedded into a belief-function framework and the missing parts can be represented as ignorance rather than filled with unjustified precise probabilities. In this thesis, the complete Bayesian network is therefore not the final object of interest. It is used as a reference structure and validation point for building and evaluating the belief-network implementation.

The aim of this thesis is practical and explanatory. It asks when an industrial user should consider a belief network, how such a model can be built, and what limitations should be expected. The thesis does not claim that belief networks should replace Bayesian networks in all diagnostic applications. If reliable data are available for all relevant priors and conditional probability tables, Bayesian inference is usually simpler, more established, and computationally more efficient. The motivation for belief networks arises in the middle ground where the diagnostic structure is sufficiently understood to model, but the quantitative information required for a fully specified probabilistic model is incomplete.

The main research question is:

\begin{quote}
When are belief networks a suitable approach for industrial fault diagnosis under incomplete probabilistic information, and how can such a model be implemented and evaluated in practice?
\end{quote}

This question is divided into the following subquestions:

\begin{enumerate}
    \item How can a fault-diagnosis problem be represented as a belief network?
    \item How can missing or weakly justified diagnostic information be represented using belief functions?
    \item How can a belief-network model be implemented and validated using the Belief Function Machine?
    \item How do belief, plausibility, and diagnostic rankings behave as probabilistic information is progressively removed?
    \item Under what practical conditions is the belief-network approach useful for industrial users, and what are its main limitations?
\end{enumerate}

The contribution of the thesis is to develop and evaluate a compact workflow for belief-network-based diagnosis. The electrical case study is used to demonstrate how the model is built, how missing information is represented, how the Belief Function Machine is used for inference, and how the outputs should be interpreted. The complete Bayesian reference model provides a consistency check under complete information, while the main analysis focuses on the behaviour of the belief network under progressively incomplete information. Section~\ref{sec:background} introduces the required background. Section~\ref{sec:case_study_diagnostic_model} presents the diagnostic model. Section~\ref{sec:building_belief_network} explains the implementation. Section~\ref{sec:experimental_results} combines the experimental design and results. Section~\ref{sec:practical_guidance} discusses when the approach should and should not be recommended. Section~\ref{sec:conclusion} concludes the thesis.

\section{Background}
\label{sec:background}

\subsection{Fault Diagnosis under Incomplete Information}
\label{sec:fault_diagnosis_uncertainty}

Fault diagnosis is the task of identifying the cause of abnormal system behaviour from available observations. In engineered systems, this usually means reasoning from symptoms to hidden faults. Symptoms may include sensor values, alarms, indicator lights, user reports, or measured signals. Faults are often not directly observable, and several different faults may produce similar symptoms. A diagnostic method must therefore manage uncertainty about the system state, the reliability of the available evidence, and the completeness of the diagnostic knowledge \cite{Venkatasubramanian2003PartI,Venkatasubramanian2003PartII,Venkatasubramanian2003PartIII,Isermann2006}.

Fault diagnosis methods are commonly classified into model-based, signal-based, data-driven, and hybrid approaches. Model-based approaches use explicit knowledge of the system structure or physics. Signal-based approaches focus on patterns in measured signals. Data-driven approaches learn diagnostic relationships from historical data. Hybrid approaches combine several sources of information because no single source is complete in many realistic applications \cite{Venkatasubramanian2003PartI,Venkatasubramanian2003PartII,Venkatasubramanian2003PartIII,GaoCecatiDing2015PartII}.

This thesis focuses on model-based graphical diagnosis under incomplete information. The electrical system used in the case study is simple enough for its structure to be described explicitly: voltage should propagate from the power supply through the control modules to the load, and certain faults interrupt this propagation. However, knowing the structure of the system does not automatically provide all probabilities needed for probabilistic diagnosis. This distinction is central to the thesis. The qualitative diagnostic model may be available, while the quantitative information may be incomplete.

Uncertainty in fault diagnosis can arise in several ways. Observations may be incomplete, for example when a technician observes that the final lamp is off but has no measurements for intermediate modules. Observations may also be ambiguous: a downstream indicator being off may be caused by a local fault or by an upstream interruption. Diagnostic knowledge may be incomplete when experts know that a relation exists between a fault and a symptom but cannot quantify that relation for every possible condition. Finally, evidence sources may conflict. These characteristics make fault diagnosis a suitable domain for studying belief networks, because the diagnostic problem is not only to choose a fault, but also to communicate when the available evidence is insufficient for a precise conclusion.

\subsection{Bayesian Networks as a Reference Framework}
\label{sec:bayesian_networks_reference}

A Bayesian network is a directed acyclic graph together with a set of probability distributions. The nodes represent variables, and the edges represent direct probabilistic dependencies. A root node is assigned a prior probability distribution. A non-root node is assigned a conditional probability table specifying the probability of its states given the states of its parents. Together, the graph and the local probability tables define a joint probability distribution over all variables \cite{Pearl1988,Heckerman1995}.

If the network variables are denoted by $X_1, \ldots, X_n$, and if $Pa_i$ denotes the parents of variable $X_i$, the joint probability distribution factorises as

\begin{equation}
P(x_1, \ldots, x_n) = \prod_{i=1}^{n} P(x_i \mid pa_i).
\label{eq:bn_factorisation}
\end{equation}

This factorisation is attractive for diagnosis because the model can be built modularly. Fault variables influence internal state variables, and internal state variables influence observations. Once evidence is observed, Bayesian inference updates the posterior probabilities of the fault variables. In the electrical case study, for example, a battery fault can affect voltage at the power supply output, which can affect downstream modules and ultimately whether the final lamp is on.

The strengths of Bayesian networks are important. They are interpretable, modular, mathematically coherent, and supported by mature inference algorithms. They can combine prior knowledge with observations, and they have been used in diagnostic applications where hidden system states must be inferred from imperfect evidence \cite{LernerParrKollerBiswas2000,VerbertBabuskaDeSchutter2017}. For this reason, the Bayesian network in this thesis is not treated as an obsolete or inappropriate method. It is used as a reference framework under complete probabilistic information.

The limitation relevant here is the requirement for precise probabilistic parameters. A complete Bayesian network requires all priors and conditional probability table entries to be specified. For larger models, especially when variables have several parents, this can require many numerical values. Some values may not be available because rare faults do not appear often enough in data, experts cannot confidently assign exact probabilities, or some parent-state combinations are physically possible but unobserved. Filling such gaps with assumptions makes inference possible, but the resulting precise posterior may suggest a level of confidence that the model inputs do not justify \cite{Shenoy2023IncompleteBN}.

\subsection{Dempster--Shafer Theory and Belief Functions}
\label{sec:dempster_shafer_belief_functions}

Dempster--Shafer theory, also known as belief-function theory or evidence theory, provides a framework for representing uncertain and incomplete information \cite{Dempster1967,Shafer1976}. Its main difference from classical probability theory is that support can be assigned to sets of hypotheses rather than only to individual hypotheses.

Let $\Theta$ be a finite frame of discernment, meaning the set of mutually exclusive and exhaustive hypotheses considered by the model. A basic probability assignment, also called a mass function, is a function

\begin{equation}
m : 2^{\Theta} \rightarrow [0,1]
\label{eq:mass_function}
\end{equation}

such that

\begin{equation}
m(\emptyset) = 0, \qquad \sum_{A \subseteq \Theta} m(A) = 1.
\label{eq:mass_constraints}
\end{equation}

The subsets $A \subseteq \Theta$ for which $m(A)>0$ are called focal elements. A Bayesian probability distribution is a special case in which all focal elements are singletons. A vacuous mass function, defined by $m(\Theta)=1$, represents complete ignorance over the frame \cite{Shafer1976,SmetsKennes1994}.

A simple diagnostic example clarifies the idea. Suppose the possible faults are battery fault, cable fault, and lamp fault. A Bayesian probability distribution must assign probability mass separately to each fault. A belief function can instead assign mass to the set \{battery fault, cable fault\} when the evidence supports an upstream electrical interruption but does not justify distinguishing between those two possibilities. This does not mean that the model splits the support equally between the two faults. It means that the support is committed to the set as a whole and not to either individual fault.

From a mass function, two important quantities are derived: belief and plausibility. For a proposition $A \subseteq \Theta$, they are defined as

\begin{equation}
Bel(A) = \sum_{B \subseteq A} m(B),
\label{eq:belief}
\end{equation}

\begin{equation}
Pl(A) = \sum_{B \cap A \neq \emptyset} m(B) = 1 - Bel(\bar{A}).
\label{eq:plausibility}
\end{equation}

Belief is the support fully committed to $A$. Plausibility is the support compatible with $A$ \cite{Shafer1976,SmetsKennes1994}. In this thesis, belief and plausibility are treated primarily as measures of committed and compatible support, not as ordinary posterior probabilities. In other interpretations of belief functions, the pair $[Bel(A),Pl(A)]$ can be related to lower and upper probability bounds, but that interpretation is not the focus of this thesis. The relevant diagnostic point is that a low belief and high plausibility for a fault means that the fault is not strongly supported individually, but remains compatible with the available evidence.

Evidence from different sources can be combined using Dempster's rule of combination:

\begin{equation}
(m_1 \oplus m_2)(A) =
\frac{\sum_{B \cap C = A} m_1(B)m_2(C)}
{1 - \sum_{B \cap C = \emptyset} m_1(B)m_2(C)}, \quad A \neq \emptyset.
\label{eq:dempsters_rule}
\end{equation}

The intuitive idea is that two distinct bodies of evidence are intersected: support is strengthened where sources agree, while conflict appears where sources support incompatible hypotheses \cite{Shafer1976,ShaferShenoyMellouli1987}. However, the rule must be used carefully. It assumes that sources are distinct and should not be combined repeatedly if they are based on the same underlying information. In diagnostic systems, this matters because several observations may originate from the same sensor chain, physical dependency, or expert judgement.

\subsection{Belief Networks and the BFM Approach}
\label{sec:belief_networks_bfm_approach}

Belief functions can be organised in graphical models. In this thesis, a belief network is understood as a graphical diagnostic model whose local uncertainty relations are represented using belief functions rather than only precise conditional probabilities. This preserves a network-like structure while allowing incomplete local information to be represented explicitly.

The work most relevant to this thesis is Shenoy's approach to inference in incomplete Bayesian networks \cite{Shenoy2023IncompleteBN}. The problem is directly aligned with the present project: what should be done when the structure of a Bayesian-style diagnostic model is known, but some priors or conditional probabilities are missing? In ordinary Bayesian inference, the network must be completed before inference can be performed. Shenoy's approach instead embeds the available network information into a belief-function model. Known probabilities can be represented as belief functions, while missing information can be represented vacuously or left partially specified.

This gives the Bayesian network a specific role in the thesis. It is not the method being promoted for incomplete cases; rather, it provides a familiar reference structure and a validation point. When all probabilistic information is available, the corresponding belief-network representation should reproduce the same probabilistic behaviour. Once this correspondence is checked, information can be removed deliberately to study how the belief network behaves under incomplete knowledge.

The implementation uses the Belief Function Machine (BFM), a MATLAB environment for reasoning with belief functions \cite{GiangShenoy2003BFM}. BFM allows a user to define variables, relations, conditional relations, valuations, and conditional valuations in the Unified Input Language. Conditional beliefs can then be converted into unconditional belief functions using Smets' conditional embedding, after which inference is performed by combining and marginalizing belief functions. In this thesis, Shenoy's incomplete-network framework provides the theoretical bridge, while BFM provides the computational environment for constructing and testing the model.

\section{Case Study and Diagnostic Model}
\label{sec:case_study_diagnostic_model}

This section introduces the diagnostic case study used to demonstrate and evaluate the method. The purpose is not to provide a full engineering specification of the electrical system or a complete listing of all probability tables. Those details belong in the appendix. The main text gives only the information needed to understand the belief-network construction and the experiments.

\subsection{Electrical System and Diagnostic Task}
\label{sec:electrical_system_diagnostic_task}

The case study is based on a modular electrical system whose overall function is to operate a load, represented by a lamp. The system consists of ten modules connected in series: one power supply module, eight control modules, and one load module. The power supply module contains a 12V lithium battery and an indicator LED. Each control module contains a switch and a switch indicator LED. The load module contains the main lamp and a lamp indicator.

The diagnostic task is to infer which hidden fault, or group of faults, explains the observed behaviour of the system. The model distinguishes between hidden fault variables, intermediate voltage variables, and observable variables. Fault variables represent possible root causes, including a power supply short, battery exhaustion, switch detachment faults, inter-module cable faults, and lamp failure. Voltage variables represent propagation of electrical power through the system, with $V_0$ denoting voltage at the power supply output and $V_1$ to $V_8$ denoting post-switch voltage after each control module. Observable variables include the power supply LED, switch indicator LEDs, the lamp indicator, the final lamp, and optional multimeter readings.

The system is intentionally simple, but it captures properties that are important in real diagnostic settings. Faults can propagate through the chain, so an upstream fault can cause downstream indicators and the final lamp to turn off. Different faults can produce similar observations, such as a detached switch, broken cable, or missing upstream voltage. Some observations may also be unavailable, for example when indicator LEDs are removed or no multimeter is available.


\subsection{Bayesian Network Reference Model}
\label{sec:bayesian_network_reference_model}

A Bayesian network reference model was developed collaboratively as part of the group project and is used here as a reference structure for the belief-network implementation. The model was built from the physical logic of the ten-module electrical system: fault nodes influence voltage propagation, voltage states influence indicator and lamp observations, and optional multimeter readings provide additional evidence. The full Bayesian network specification, including the graph structure, priors, conditional probability tables, diagnostic scenarios, and implementation files, is documented in the shared project repository \cite{BNRepository2026}.

The numerical probabilities in the Bayesian network should not be interpreted as real industrial reliability data. Most conditional relations are deterministic because they follow the circuit logic. For the fault priors and the few non-deterministic relations, ChatGPT-5.5 was used as a simulated domain expert to propose plausible and internally consistent values. The purpose of these values is not to estimate true field failure rates, but to create a realistic toy-domain reference model that produces meaningful diagnostic ambiguity. This allows the thesis to focus on its main question: when and how belief networks can be useful when such probabilistic information is incomplete, unavailable, or weakly justified.

The reference model has two purposes. First, it provides a complete and interpretable diagnostic model under full probabilistic information. Second, it provides the structure from which the belief-network implementation is derived.
\subsection{Diagnostic Scenarios and Scope}
\label{sec:diagnostic_scenarios_scope}

The diagnostic scenarios are used to evaluate model behaviour under different evidence patterns. Some scenarios are relatively clear, such as a lamp fault where upstream indicators function but the final lamp is off. Other scenarios are more ambiguous, such as an upstream interruption where several downstream indicators are off. Some scenarios remove observations, for example by assuming that control-module LEDs are unavailable or that no multimeter measurement can be taken. The scenarios do not define the model structure; they are used to test the behaviour of a model designed from the physical system.

Several simplifying assumptions are made to keep the project manageable. The model does not include all possible electrical faults. Diode faults, LED reliability faults, internal wiring faults, and detailed polarity faults are outside the selected scope. The system is also discretised: voltage is represented using states such as 12V and 0V, and observations are represented as on or off. This abstraction is appropriate because the objective is not electrical simulation, but the representation of diagnostic uncertainty under incomplete probabilistic information.

\section{Building the Belief Network}
\label{sec:building_belief_network}

This section describes the implementation workflow. Its purpose is partly tutorial: it explains how a practitioner could move from a known diagnostic structure to a belief-network representation, and how missing information can be preserved rather than artificially completed.

\subsection{Representing Variables, Relations, and Evidence}
\label{sec:variables_relations_evidence}

The first step is to define the variables and states of the diagnostic model. Fault variables represent possible root causes, voltage variables represent internal propagation states, and observable variables represent LED states, lamp states, and measurements. In the belief-network representation, these variables are defined with discrete frames. For example, a battery variable may have states corresponding to functioning or exhausted, while a voltage variable may have states corresponding to 12V or 0V.

The second step is to define relations between variables. Deterministic physical relations, such as voltage propagation through a connected cable and closed switch, can be represented as belief valuations that assign mass to compatible configurations. Conditional diagnostic relations can be represented using conditional valuations. When the relation is fully known and probabilistic, the valuation can behave like a Bayesian conditional table. When the relation is only partially known, the valuation can leave mass on broader sets of possibilities.

Evidence is represented by assigning support to observed states. For example, observing that a given LED is off creates evidence for the corresponding observable variable. Evidence does not directly identify the root fault; instead, it is combined with the model's structural relations and propagated through the network to obtain beliefs and plausibilities for fault variables.

\subsection{Representing Missing Information}
\label{sec:representing_missing_information}

The central implementation problem is how to represent missing probabilities without replacing them with unjustified precise values. In a Bayesian network, a missing prior or missing conditional probability row prevents standard inference unless the model is completed. In a belief network, missing information can be represented as ignorance.

For example, if the prior probability of a fault variable is unknown, the belief network can assign mass to the full frame of that variable. This states that the model does not currently have support favouring one state over another. Similarly, if a conditional probability row is missing for a parent configuration, the belief-network representation can avoid inventing a precise distribution for that row. The model does not create missing knowledge; it preserves the fact that the knowledge is missing and propagates that incompleteness to the diagnostic output.

This treatment is the main reason belief networks may be useful for industrial users. In early-stage systems, rare-fault settings, or expert-supported diagnosis, the structure of a system may be known before reliable probability estimates are available. A belief network allows the modeller to use the available structural knowledge while keeping incomplete quantitative knowledge visible.

\subsection{Implementation with the Belief Function Machine}
\label{sec:implementation_bfm}

The implementation uses the Belief Function Machine in MATLAB \cite{GiangShenoy2003BFM}. BFM stores knowledge as belief functions and uses a text-based Unified Input Language. The workflow used in this thesis follows the standard BFM process. First, variables, relations, conditional relations, valuations, and conditional valuations are written in a UIL file. Second, the file is translated into BFM format using \texttt{uil2bm}. Third, the generated model is loaded in MATLAB. Fourth, conditional beliefs are converted into unconditional beliefs using \texttt{condiembed}. Fifth, the relevant beliefs are retained with \texttt{keepbel}. Finally, inference is performed using \texttt{solve} and displayed with \texttt{showbel}.

The important conceptual step is Smets' conditional embedding. This converts conditional belief valuations into unconditional belief functions on the joint domain of the child and conditioning variables. This is what allows Bayesian-style conditional knowledge to be represented inside the belief-function framework. In the present thesis, BFM is therefore not treated as a production-ready industrial platform. It is a research implementation that makes it possible to construct and test the belief-network approach.

\subsection{Worked Diagnostic Example}
\label{sec:worked_diagnostic_example}

Consider a simplified diagnostic situation in which the final lamp is off and several downstream indicators are also off. In a complete Bayesian network, the model would update the posterior probabilities of possible faults such as battery failure, cable fault, switch detachment, or lamp failure. The result would be a precise probability for each fault.

In a belief network under incomplete information, the result may be different. If the model has enough evidence to support an upstream interruption but lacks enough information to distinguish between battery failure and an early cable fault, it can assign support to a set of possible faults. The belief in each individual fault may remain low, while the plausibility of several upstream faults remains high. This tells the technician that the evidence points to a region of the system, but that the model does not have enough information to isolate a single component. The usefulness of the belief network therefore lies not in always selecting one fault with high confidence, but in communicating when the available information is insufficient for precise localization.

\section{Experimental Design and Results}
\label{sec:experimental_results}

This section evaluates the belief-network implementation developed in Section~\ref{sec:building_belief_network}. The purpose is twofold. First, the complete belief-function model is validated against the complete Bayesian network reference model to show that the implementation behaves correctly when all probabilistic information is available. Second, the model is tested under progressively incomplete information to examine whether belief-network inference remains diagnostically useful when priors, measurements, or observation relations are missing or weakened.

The experiments therefore serve both a technical and a practical role. Technically, they check that the BFM model correctly represents the diagnostic structure. Practically, they show how the model behaves in situations that resemble industrial diagnosis, where the system structure may be understood but some probability estimates or observations are not fully reliable. The aim is not to show that belief networks are more accurate than Bayesian networks under complete information. Rather, the aim is to determine whether the belief-network output remains useful when precise probabilistic modelling is no longer justified.

\begin{table}[H]
\caption{Experimental levels used in the belief-network evaluation.}\label{tab:experiment_levels}
\centering
\small
\begin{tabularx}{\textwidth}{p{1.2cm}p{3.1cm}X}
\toprule
Level & Information condition & Purpose \\
\midrule
L0 & Complete BFM model & Validate that the belief-network implementation reproduces the complete BN under precise information. \\
L1 & Local prior ignorance & Replace only the locally relevant fault priors by ignorance. \\
L2 & Fault-family prior ignorance & Replace broader relevant prior families, such as switch/cable priors or power-subsystem priors, by ignorance. \\
L4 & Combined incompleteness & Combine missing measurements, missing prior knowledge, and weakened observation reliability. \\
\bottomrule
\end{tabularx}
\end{table}

The results are interpreted using five criteria. Belief measures direct committed support for a fault state, while plausibility measures whether that fault remains compatible with the evidence. The width between belief and plausibility is used as an indicator of unresolved evidential support. Diagnostic ranking is used to check whether the relevant fault or fault family remains among the main candidates. Finally, the empty-set mass is reported separately as a measure of conflict between the combined evidence and the model assumptions.

\subsection{Complete-information validation: L0}
\label{sec:l0_validation}

The L0 model was used as a validation bridge between the complete BN and the BFM implementation. In this condition, all priors and diagnostic relations from the BN were encoded in the BFM model without introducing additional ignorance. As expected, belief and plausibility collapsed to the same value for all queried singleton fault states. This means that the L0 belief-function model behaves like a precise probabilistic model and can be compared directly with the BN posteriors.

\begin{table}[H]
\caption{L0 validation results under complete information.}\label{tab:l0_results}
\centering
\small
\begin{tabularx}{\textwidth}{p{0.9cm}p{2.7cm}p{4.0cm}X}
\toprule
Scen. & Diagnostic case & Main L0 output & Interpretation \\
\midrule
7 & Battery exhaustion & \texttt{Batt = exhausted}: Bel.=Pl.=1.00 & BFM reproduces the complete BN conclusion. \\
8 & Module-3 interruption & \texttt{SW3 = detached}: 0.607; \texttt{C3 = broken}: 0.405 & BFM reproduces the BN ranking and localizes the fault around module 3. \\
11 & Sparse downstream failure & Support distributed across several switch/cable faults & Sparse evidence already prevents precise localization under complete information. \\
14 & PSU short with battery discharge & \texttt{PSF = yes}: 1.00; \texttt{Batt = exhausted}: 1.00 & BFM reproduces the compound upstream diagnosis. \\
\bottomrule
\end{tabularx}
\end{table}

The validation confirms that, when all information is precise, the BFM is a faithful translation of the BN for the tested cases. Scenario 8 is particularly important because it shows that the model can reproduce not only decisive diagnoses, but also ambiguous rankings around a fault transition. The additional value of the belief-network representation is therefore not expected at L0, but in the later levels where information is deliberately made incomplete.

\subsection{Local prior ignorance: L1}
\label{sec:l1_local_prior_ignorance}

The L1 condition tested local prior ignorance. Only the scenario-specific prior information was replaced by ignorance, while the diagnostic structure, observations, and deterministic relations remained available. This represents a realistic situation in which a company may understand the system structure but no longer trust the failure-rate estimate for a specific component, for example after a supplier change or redesign.

\begin{table}[H]
\caption{L1 results under local prior ignorance.}\label{tab:l1_results}
\centering
\small
\begin{tabularx}{\textwidth}{p{0.9cm}p{3.0cm}p{3.9cm}X}
\toprule
Scen. & Missing information & Main L1 output & Interpretation \\
\midrule
7 & Battery-related local prior & \texttt{Batt}: Bel.=Pl.=1.00 & Direct battery evidence is strong enough to preserve the diagnosis. \\
8 & Local priors for \texttt{SW3} and \texttt{C3} & \texttt{SW3}, \texttt{C3}: Bel.=0, Pl.=1.00 & The model preserves both local explanations but refuses to rank them precisely. \\
11 & Local priors for \texttt{SW6} and \texttt{C6} & \texttt{SW6}, \texttt{C6}: Bel.=0, Pl.=1.00 & Relevant local faults remain possible, but sparse evidence prevents precise commitment. \\
14 & PSU-short and battery-related local priors & \texttt{PSF}, \texttt{Batt}: Bel.=Pl.=1.00 & Direct measurements dominate missing local prior information. \\
\bottomrule
\end{tabularx}
\end{table}

The most informative L1 result occurred in Scenario 8. Under complete information, the BN and L0 BFM preferred \texttt{SW3 = detached} over \texttt{C3 = broken}. This precise preference was mainly driven by the relative priors of switch and cable faults. When those local priors were removed, the BFM no longer had a justified basis for choosing between the two. Consequently, both \texttt{SW3} and \texttt{C3} received belief 0 and plausibility 1.00. This should not be interpreted as failure. The model still localized the fault to the module-3 transition, but it refused to make an unjustified subcomponent-level distinction.

Scenario 11 showed the same behaviour under weaker observations. The complete model already produced a broad distribution because the available evidence did not identify the exact interruption point. After the local module-6 priors were replaced by ignorance, \texttt{SW6} and \texttt{C6} remained fully plausible but received no committed belief. This is a conservative diagnostic output: the model keeps the relevant local hypotheses open while making clear that the evidence is not strong enough for precise ranking.

\subsection{Fault-family prior ignorance: L2}
\label{sec:l2_fault_family_prior_ignorance}

The L2 condition removed broader prior families. In the control-chain cases, all switch and cable priors were replaced by ignorance; in the upstream cases, power-subsystem priors were replaced by ignorance. This simulates a more severe but realistic setting in which the system structure is known, but reliability estimates for a whole component family are unavailable or no longer trusted.

\begin{table}[H]
\caption{L2 results under fault-family prior ignorance.}\label{tab:l2_results}
\centering
\small
\begin{tabularx}{\textwidth}{p{0.9cm}p{3.0cm}p{4.0cm}X}
\toprule
Scen. & Missing information & Main L2 output & Interpretation \\
\midrule
7 & Upstream power-family priors & \texttt{Batt}: Bel.=Pl.=1.00 & Direct evidence still identifies battery exhaustion. \\
8 & All switch/cable priors & Downstream switch/cable faults remain plausible; upstream faults ruled out & Diagnosis broadens, but remains physically constrained to the downstream region from module 3. \\
11 & All switch/cable priors with sparse evidence & Broad plausibility over switch/cable candidates & The model becomes cautious because neither observations nor priors justify exact localization. \\
14 & Upstream power-family priors & \texttt{PSF}, \texttt{Batt}: Bel.=Pl.=1.00 & Direct PSU-short and battery measurements preserve the compound diagnosis. \\
\bottomrule
\end{tabularx}
\end{table}

Compared with L1, L2 widened the output, but the uncertainty remained structured. In Scenario 8, faults before module 3 were still ruled out because the evidence showed that power reached modules 1 and 2. The plausible region expanded only among downstream switch and cable faults compatible with the observed loss of signal. Thus, the model did not become randomly ignorant; it preserved the physical constraints imposed by the remaining evidence. In Scenarios 7 and 14, where direct measurements identified the upstream faults, L2 did not reduce diagnostic specificity: battery exhaustion and PSU short were still assigned belief/plausibility 1.00 where appropriate.

\subsection{Combined incompleteness: L4}
\label{sec:l4_combined_incompleteness}

The L4 condition tested more realistic mixed incompleteness. Two derived scenarios were used. Scenario 15 was derived from the battery-discharge cases by removing the PSU-short measurement. The battery measurement still indicated 0V, but the model no longer had the discriminating observation needed to confirm or reject a hidden PSU short. Scenario 16 was derived from the module-3 interruption case by removing the key boundary indicator \texttt{I3}; the model still knew that power reached modules 1 and 2 and that later indicators were off, but the first point of failure was less directly observed.

\begin{table}[H]
\caption{L4 results under combined incompleteness.}\label{tab:l4_results}
\centering
\small
\begin{tabularx}{\textwidth}{p{0.9cm}p{3.2cm}p{4.0cm}X}
\toprule
Scen. & Combined incompleteness & Main L4 output & Interpretation \\
\midrule
15 & PSU-short measurement unavailable; power priors ignorant; visual rules weakened & \texttt{Batt}: Bel.=Pl.=1.00; \texttt{PSF}: Bel.=0, Pl.=1.00 & Battery exhaustion is directly supported, while hidden PSU short remains unresolved but plausible. \\
16 & Boundary indicator \texttt{I3} unavailable; switch/cable priors ignorant; visual rules weakened & Downstream \texttt{SW3--SW8} and \texttt{C3--C9}: Bel.=0, Pl.=1.00 & Exact localization is lost, but the model still rules out upstream causes and preserves a useful downstream diagnostic region. \\
\bottomrule
\end{tabularx}
\end{table}

The L4 results show the most realistic form of belief-function behaviour. In Scenario 15, the model directly supports battery exhaustion, but it does not rule out a hidden PSU short as a possible upstream cause. A technician would therefore be warned not to assume a battery-only failure without checking the PSU short. In Scenario 16, combined missing priors, weakened visual rules, and the absence of the boundary indicator prevented precise component ranking. However, the BFM still ruled out upstream faults such as battery exhaustion, PSU short, and faults before module 3. This result illustrates a useful degradation pattern: the model loses exact specificity but preserves a physically meaningful diagnostic region.

\subsection{Conflict and interpretation}
\label{sec:conflict_interpretation}

The empty-set mass, $m(\emptyset)$, was retained as a separate indicator of conflict between the combined evidence and the model assumptions. In L0, conflict was high because precise priors encoded faults as rare while the evidence in each scenario indicated that a fault had occurred. The BN normalizes this tension away when conditioning on evidence, while the BFM exposes it as conflict.

\begin{table}[H]
\caption{Empty-set mass across experimental levels.}\label{tab:conflict_results}
\centering
\small
\begin{tabular}{ccccc}
\toprule
Scenario & L0 & L1 & L2 & L4 \\
\midrule
7  & 0.920 & 0.005 & 0.000 & -- \\
8  & 0.959 & 0.173 & 0.085 & -- \\
11 & 0.683 & 0.084 & 0.084 & -- \\
14 & 0.995 & 0.000 & 0.000 & -- \\
15 & --    & --    & --    & 0.000 \\
16 & --    & --    & --    & 0.085 \\
\bottomrule
\end{tabular}
\end{table}

The reduction in conflict from L0 to L1 and L2 suggests that part of the conflict in the complete model was caused by overly precise prior assumptions clashing with strong fault evidence. Replacing unsupported priors with ignorance reduced this artificial prior-evidence tension. This is relevant for industrial use: if reliability estimates are weak or outdated, forcing them into a precise BN may create confident but poorly justified rankings. The belief-network approach instead makes the loss of knowledge visible through wider belief-plausibility gaps and, in some cases, lower conflict.

Overall, the experiments show three main patterns. First, when evidence is decisive, as in Scenarios 7 and 14, the BFM can still identify the correct faults even when prior information is weakened. Second, when the evidence localizes a fault region but does not justify a precise component distinction, as in Scenario 8, the BFM preserves the relevant candidates as plausible rather than inventing a precise ranking. Third, when both observations and prior knowledge are weak, as in Scenarios 11 and 16, the model becomes broad but still physically constrained. These findings support the practical interpretation of belief networks: they are not mainly useful because they are more precise than Bayesian networks, but because they can show when precision is no longer justified.
\section{Practical Applicability and Limitations}
\label{sec:practical_guidance}

\subsection{When Belief Networks Are Recommended}
\label{sec:when_recommended}

Belief networks are recommended for diagnostic settings where the structure of the system is known well enough to model, but precise probabilities are incomplete, weakly justified, or unavailable. This can occur in early-stage product development, rare-fault diagnosis, safety-critical systems with limited failure data, or industrial settings where expert knowledge is qualitative rather than numerical.

In these cases, the advantage of a belief network is not that it automatically gives a more accurate diagnosis. Its advantage is that it can reason with the information that is available while preserving the fact that some information is missing. This is valuable when a precise posterior probability would be misleading because it depends on arbitrary completion assumptions. A belief network is especially useful when the practitioner wants to know not only which fault is supported, but also whether alternative explanations have actually been ruled out.

\subsection{When Belief Networks Are Not Recommended}
\label{sec:when_not_recommended}

Belief networks are not always the best choice. If sufficient data or reliable expert estimates are available for the relevant priors and conditional probability tables, a Bayesian network is likely to be simpler, faster, easier to communicate, and better supported by software. Similarly, if the end user requires a single posterior probability or a fully automated ranking without explanation of ignorance, the additional interpretive complexity of belief and plausibility may not be worthwhile. This is consistent with comparative work on Bayesian and Dempster--Shafer reasoning in fault diagnosis, where Bayesian reasoning is often preferable when information is sufficiently complete and non-conflicting, while Dempster--Shafer reasoning becomes more informative when incompleteness or conflict is central to the diagnostic problem \cite{VerbertBabuskaDeSchutter2017}.

Belief networks are also risky when evidence sources are not distinct. Dempster's rule assumes that combined bodies of evidence are distinct. If the same underlying information is counted repeatedly, the model can become overconfident for the wrong reasons. This is important in diagnostic applications where several observations may be generated by the same physical dependency or sensor chain.

\subsection{Scalability, Tooling, and Technology Readiness}
\label{sec:scalability_tooling_readiness}

The main technical limitation of belief networks is computational complexity. Because belief functions can assign mass to sets of hypotheses, the number of focal elements can grow quickly. This can make inference more expensive than ordinary Bayesian inference, especially in large networks with many variables and states. Practical deployment would require careful modeling choices, modularisation, local computation, and possibly approximation methods \cite{ShaferShenoyMellouli1987,Shenoy2023IncompleteBN}.

Software maturity is another limitation. Bayesian networks benefit from a larger ecosystem of tools, libraries, tutorials, and industrial examples. Belief-network tooling is more specialised. The Belief Function Machine provides a useful research environment, but it is not as widely adopted or as easy to integrate into standard industrial software workflows as common Bayesian-network tools. MATLAB should therefore be understood as the research platform used in this thesis, not as the only possible implementation environment for the underlying method.

From a technology-readiness perspective, belief networks should be presented carefully. The underlying theory is mature, and research implementations exist, but the path from a research-grade implementation to industrial deployment is not automatic. The method requires expertise in uncertainty modeling, careful treatment of evidence dependence, validation on realistic diagnostic tasks, and clear guidance for interpreting belief and plausibility outputs.

\subsection{Future Work}
\label{sec:future_work}

Future work should test the approach on larger and more realistic diagnostic systems, compare different ways of representing partial conditional knowledge, and investigate approximation strategies for larger belief networks. More work is also needed on user-facing diagnostic interfaces: industrial users would need outputs that explain not only which faults are plausible, but also why the model remains uncertain and which additional measurements would reduce that uncertainty. Finally, future implementations could explore Python, Java, or C++ environments to separate the practical method from the MATLAB-based research prototype used here.

\section{Conclusion}
\label{sec:conclusion}

This thesis investigates belief networks as a practical approach for industrial fault diagnosis under incomplete probabilistic information. The central motivation is that diagnostic structure may be available even when the exact probabilities required for a fully specified Bayesian network are not. In such cases, completing the model with artificial point estimates can produce precise results that hide the incompleteness of the underlying knowledge.

The thesis uses a modular electrical system as a case study. A complete Bayesian network is used as a reference structure and validation point, while the main focus is the construction and evaluation of a belief-network implementation using the Belief Function Machine. The experiments progressively remove probabilistic information in order to study how belief, plausibility, and diagnostic rankings behave under increasing incompleteness.

The contribution is a practical recommendation framework. Belief networks should be considered when the diagnostic structure is known, probabilities are incomplete or weakly justified, and preserving ignorance is important for decision support. They are less appropriate when reliable probabilities are available, when computational scalability is critical, or when users require simple single-probability outputs. belief networks are not a universal replacement for Bayesian networks, but they can be a useful expert-supported diagnostic approach when missing probabilistic information is itself an important part of the problem.

\appendix

\section{Additional Dataset Details}
\label{app:additional_dataset_details}

[TO FILL: Add additional system details, scenario tables, and dataset construction notes that are too detailed for the main text.]

\section{Additional Bayesian Network Details}
\label{app:additional_bn_details}

[TO FILL: Add the full directed acyclic graph specification, priors, conditional probability tables, and validation checks.]

\section{Additional Belief Network Details}
\label{app:additional_belief_network_details}

[TO FILL: Add detailed belief-network construction notes, mass functions, BFM implementation details, and any equations that are too technical for the main text.]

\section{Additional Results}
\label{app:additional_results}

[TO FILL: Add additional scenario outputs, extended tables, and robustness checks.]


% ---- Bibliography ----
\bibliographystyle{splncs04}
\bibliography{references_amended_samplepaper}

\end{document}
